\documentclass[11pt,a4paper]{article}
\usepackage{nesarticle}
\begin{document}
\heading{Homework 1 · Instructor solutions}{Lectures 1 and 2 · Parts A and B}
\paragraph{Assessment.} Part A is a 120-minute formative mock, marked out of 100 for diagnosis only: mistakes do not reduce the final course grade. Part B is the ordinary assessed homework, redesigned as an instructional training set, also out of 100. Do not add or average the two scores. The three assessed homeworks remain equally weighted within the 20\% homework component. Both parts are handwritten. No deadline is invented here; the revised release schedule remains blocks 3/5/7.
\paragraph{Timing.} The selected Part A is demanding for 120 minutes: ten claims, two multi-part problems and two three-sentence answers. Treat completion and where a student runs out of time as diagnostic information. Do not require a textbook-length proof when a concise valid argument suffices. Part B has no time limit and is designed for deliberate practice with notes.
\paragraph{Part A marking.} Ten true/false claims: 4 points each, normally 1 for the verdict and 3 for a valid justification; a bare unsupported label earns zero as announced. Problem A2 subparts (a)--(e): 3, 5, 6, 3, 3 points. Problem A3 subparts (a)--(e): 4, 6, 12, 4, 4 points. Short answers: 5 each. Give partial credit for valid intermediate reasoning and accept equivalent counterexamples.
\paragraph{True/false key.} F, T, T, F, F; T, F, F, F, T.
\section*{A1. True/false solutions}
\clearpage\task{1. There is always a winner}{False}
The relation is complete: for distinct $x,y$, some coordinate differs, and the bundle with the larger value in that coordinate is weakly preferred under the rule. Both comparisons can hold.

It is not transitive. Let $x=(1,0)$, $y=(0,1)$ and $z=(2,0)$. Then $x\pref y$ because $x_1>y_1$, and $y\pref z$ because $y_2>z_2$. But $x\ne z$, $x_1<z_1$ and $x_2=z_2$, so $x\not\pref z$. Thus the relation is not rational. A strict cycle is not required: a failure of weak transitivity suffices.
\task{2. A jump in the score}{True}
Set $s(x)=x_1+x_2$ and $f(t)=t+\mathbf1\{t\ge1\}$. This $f$ is strictly increasing on $[0,\infty)$: within each piece it has slope 1, and crossing 1 adds a positive jump. Consequently
\[
 u(x)\ge u(y)\quad\Longleftrightarrow\quad s(x)\ge s(y).
\]
The continuous function $s$ represents the same preferences. More explicitly, if $x^n\pref y^n$ and the bundles converge, $s(x^n)\ge s(y^n)$ passes to $s(x)\ge s(y)$, hence $x\pref y$. The displayed index is discontinuous, but its represented preferences are continuous.
\task{3. Mixing without a concave index}{True}
On the strictly positive domain, $u>0$, so its logarithm is a valid strictly increasing transformation:
\[
 g(x)=\log u(x)=2\log x_1+2\log x_2.
\]
Its Hessian is diagonal with entries $-2/x_i^2<0$, hence $g$ is strictly concave on the convex domain $\R_{++}^2$. If distinct $x,y\pref z$ and $0<\theta<1$, then
\[
 g(\theta x+(1-\theta)y)>\theta g(x)+(1-\theta)g(y)\ge g(z).
\]
Thus preferences are strictly convex. The original index need not be concave; strict concavity of this alternative representation is sufficient. The domain is essential: on $\R_+^2$, distinct bundles on one axis have utility zero, as does their midpoint, violating strict convexity of preferences.
\clearpage
\task{4. Somewhere better means everywhere higher}{False}
Use $u(x)=x_1-x_2$ on $\R_+^2$. Representation gives rationality; continuity follows from the linear index. Every upper contour set is the intersection of the orthant with a half-space $x_1-x_2\ge k$, hence convex. At every bundle, an arbitrarily small increase in $x_1$ is feasible in the consumption domain and strictly improves utility, establishing LNS.

But $(1,2)\gg(0,0)$ and $u(1,2)=-1<0=u(0,0)$. LNS requires a better bundle in every neighbourhood, not improvement from every coordinatewise increase. A bad is allowed: monotonicity was not assumed.
\task{5. A hilltop or just a landing?}{False}
Take $u(x)=\max\{0,x_1-1\}$ and $x^*=(1/2,1/2)$. This function is continuous. Within Euclidean distance $1/4$ of $x^*$, $x_1<1$, hence $u(x)=0=u(x^*)$. Therefore $x^*$ is a local maximiser under the explicitly non-strict definition. It is also a local minimiser; a plateau can have both properties. Yet $u(2,1/2)=1>u(x^*)$, so it is not a global maximiser.

To check convex preferences, for each reference bundle $z$,
\[
 \{x\in\R_+^2:u(x)\ge u(z)\}=
 \begin{cases}
 \R_+^2,&z_1\le1,\\
 \{x\in\R_+^2:x_1\ge z_1\},&z_1>1.
 \end{cases}
\]
Both sets are convex. Thus the example meets every hypothesis. Quasiconcavity permits a flat region followed by improvement; concavity would make every local maximiser global. No LNS assumption was made in this question.
\task{6. A different currency, the same shopping basket}{True}
For any $t>0$, $tp\cdot x\le tm$ iff $p\cdot x\le m$, so $B(tp,tm)=B(p,m)$. The same objective on the same feasible set has the same full argmax set. Uniqueness and differentiability are unnecessary. Equality even holds if both argmax sets are empty; the question does not infer existence merely from a notation. Scaling prices alone, or income alone, is a different experiment.
\task{7. Strict curvature, a unique choice}{False}
Let $u(x)=x_1^2+x_2^2$, $p=(1,1)$ and $m=1$. The Hessian is $2I$, so the function is strictly convex and continuous. On the budget,
\[
 x_1^2+x_2^2\le(x_1+x_2)^2\le1.
\]
Equality holds at both $(1,0)$ and $(0,1)$, giving two distinct optimal bundles. Strict convexity of preferences is the ordinal sufficient condition for at most one optimum on a convex feasible set. It must not be confused with strict convexity of the function being maximised. Continuity and compactness do provide existence here.
\clearpage
\task{8. All boxes ticked}{False}
Let $u(x)=(2x_1+x_2-3)^3$, $p=(1,1)$ and $m=2$. The cube is strictly increasing on the entire real line, so preferences coincide with those represented by $2x_1+x_2$. They are convex and strongly monotone. The displayed function is a polynomial, hence continuously differentiable.

At $x^*=(1,1)$, $u(x^*)=0$ and
\[
 \nabla u(x)=3(2x_1+x_2-3)^2(2,1),\qquad \nabla u(x^*)=(0,0).
\]
Choose $\lambda=\mu_1=\mu_2=0$. Stationarity $u_i-\lambda p_i+\mu_i=0$, primal feasibility, dual feasibility and complementary slackness all hold; the budget even binds. Nevertheless $(2,0)$ is affordable and has utility 1, so $x^*$ is not optimal.

A strictly increasing transformation can have zero derivative at a point. Quasiconcavity/convex preferences alone do not make every KKT point globally optimal. Concavity of the differentiable index gives a standard sufficient condition; it is absent here.
\task{9. Left on the shelf}{False}
For $p=(1,1)$, $m=1$, use
\[
 u(x)=2\log(1+x_1)+\log(1+x_2).
\]
Partial derivatives are positive, and the Hessian is diagonal with strictly negative entries, so this utility is strongly increasing and strictly concave. At $x^*=(1,0)$,
\[
 u_1(x^*)=\frac{2}{2}=1,\qquad u_2(x^*)=\frac{1}{1}=1.
\]
The KKT conditions hold with $\lambda=1$, $\mu_1=\mu_2=0$. Concavity makes them sufficient; strict concavity gives uniqueness. Equivalently, the supporting-plane inequality gives $u(x)\le u(x^*)+\nabla u(x^*)\cdot(x-x^*)\le u(x^*)$ for every affordable $x$.

At a corner with good 1 consumed, stationarity gives $u_1/p_1=\lambda$ and $u_2/p_2=\lambda-\mu_2/p_2\le\lambda$. Equality is possible even though $x_2=0$. Binding nonnegativity does not imply a strictly positive bound multiplier.
\clearpage
\task{10. Can a smooth value hide several choices?}{True}
Continuity on the nonempty compact budget set gives an optimum; positive prices are essential for compactness here. LNS implies $p\cdot x^*=m$ for every optimum: with slack, a sufficiently small neighbourhood stays affordable and contains a strictly preferred bundle.

Fix \emph{any} optimiser $x^*$. For a price vector $q$ sufficiently close to $p$, define
\[
 F(q)=v(q,q\cdot x^*).
\]
Prices and the new wealth are positive nearby since $p\gg0$ and $p\cdot x^*=m>0$. The fixed bundle $x^*$ is feasible at $(q,q\cdot x^*)$, so $F(q)\ge u(x^*)$. At $q=p$, budget exhaustion gives $F(p)=v(p,m)=u(x^*)$. Thus $F$ has a local minimum at $p$. Differentiability and the chain rule imply
\[
 0=F_{q_i}(p)=v_{p_i}(p,m)+v_m(p,m)x_i^*.
\]
Because $v_m>0$,
\[
 x_i^*=-\frac{v_{p_i}(p,m)}{v_m(p,m)}.
\]
Every optimiser must have these same coordinates, proving uniqueness. The proof does not assume a differentiable demand function or uniqueness in advance. Merely invoking a version of Roy's identity that already assumes uniqueness would be circular; the fixed-bundle argument avoids that problem.

\clearpage\section*{Problem 1. Two tickets, one weekend}
\paragraph{(a) Rankings and rationality.} All three programmes have utility 3: the score pairs are $(6,3)$, $(3,9)$ and $(3,4)$. Thus they are mutually indifferent. A real-valued representation gives completeness and transitivity. The minimum of two continuous linear functions is continuous, so the represented weak-preference graph is closed.
\paragraph{(b) Geometry.} The level-6 indifference set is the polyline connecting
\[
 (0,6),\qquad(12/5,6/5),\qquad(6,0).
\]
The weak upper contour set is
\[
 \{x\ge0:2x_1+x_2\ge6,\quad x_1+3x_2\ge6\}.
\]
For $0\le x_1\le12/5$, the indifference set follows $x_2=6-2x_1$. Mira's score is the smaller one there (strictly smaller before the kink). For $12/5\le x_1\le6$, it follows $x_2=(6-x_1)/3$, and Anton's score is the smaller one. At the kink both scores are 6. The upper contour is above both lines, not a union of the two half-spaces.
\paragraph{(c) Properties.} At any reference utility level $k\ge0$, the upper contour is the intersection of $X$ with two half-spaces $2x_1+x_2\ge k$ and $x_1+3x_2\ge k$, so preferences are convex. They are not strictly convex: $x=(0,4)$ and $y=(1,2)$ both have utility 4, and their midpoint $(1/2,3)$ also has utility 4.

Both linear scores have strictly positive coefficients. If a bundle increases coordinatewise with at least one strict improvement, \emph{both} scores strictly rise. More precisely, both score increases have a strictly positive minimum, so their minimum also rises strictly. This proves strong monotonicity. Increasing either coordinate by an arbitrarily small positive amount then proves LNS throughout the orthant.
\paragraph{(d) Displayed scores.} Since $u(X)=[0,\infty)$, squaring is strictly increasing on the attained range and preserves every comparison. But $v$ is not concave. On $x_2=0$, $u(x)=x_1$ and $v(x)=x_1^2$. For example,
\[
 v(1,0)=1<\tfrac12v(0,0)+\tfrac12v(2,0)=2.
\]
Changing the numerical scale need not preserve concavity.
\paragraph{(e) Convex representation.} Set $x=(0,6)$ and $y=(6,0)$. Both have utility 6, while their midpoint $(3,3)$ has utility 9. Any representation $g$ must satisfy $g(x)=g(y)<g((x+y)/2)$, whereas convexity would require $g((x+y)/2)\le[g(x)+g(y)]/2=g(x)$. Thus no convex representing function exists. Convexity of preferences does not imply convexity of a utility index.
\clearpage
\section*{Problem 2. One more lesson, one less coffee}
\paragraph{(a) Existence, exhaustion and uniqueness.} The budget is nonempty and compact for $p>0,m>0$, and utility is continuous, so an optimum exists. Utility is strictly increasing in $a$ and $q$, hence an optimum exhausts the budget. Substitute $a=m-pq$ to obtain
\[
 \max_{0\le q\le m/p}\{m-pq+6\log(1+q)\}.
\]
The second derivative with respect to $q$ is $-6/(1+q)^2<0$. The reduced problem has a unique maximiser; budget exhaustion then uniquely determines $a$. The original utility is not strictly concave in the pair $(a,q)$ because it is affine in $a$; do not use that false premise to justify uniqueness.
\paragraph{(b) Kuhn--Tucker system.} With constraints written as nonnegative,
\[
 \mathcal L=a+6\log(1+q)+\lambda(m-a-pq)+\mu_a a+\mu_q q.
\]
Stationarity:
\[
 1-\lambda+\mu_a=0,\qquad
 \frac6{1+q}-\lambda p+\mu_q=0.
\]
Primal feasibility: $a,q\ge0$, $m-a-pq\ge0$. Dual feasibility: $\lambda,\mu_a,\mu_q\ge0$. Complementary slackness:
\[
 \lambda(m-a-pq)=0,\qquad \mu_a a=0,\qquad\mu_q q=0.
\]
The objective is concave and differentiable on a neighbourhood of the consumption domain; the constraints are affine. KKT candidates satisfying these conditions are globally optimal. The reduced problem in (a) already establishes uniqueness.
\paragraph{(c) Demand.} The reduced derivative is $-p+6/(1+q)$, strictly decreasing in $q$. Thus
\[
 q^*(p,m)=\min\left\{\frac m p,\ \max\left\{\frac6p-1,0\right\}\right\},\qquad
 a^*(p,m)=m-pq^*(p,m).
\]
The full classification is
\[
 (a^*,q^*)=
 \begin{cases}
 (m,0),&p\ge6,\\[1mm]
 (0,m/p),&0<p<6,\quad 0<m\le6-p,\\[1mm]
 (m-6+p,6/p-1),&0<p<6,\quad m>6-p.
 \end{cases}
\]
At $p=6$, demand for tutoring is still uniquely zero: the reduced derivative is zero only at zero and negative thereafter. At $m=6-p$ with $p<6$, $a=0$ and $q=6/p-1=m/p$; this is the boundary between the two positive-tutoring regimes, not a multiplicity of optima.

The corresponding multipliers are
\[
 \lambda=
 \begin{cases}
 1,&p\ge6,\\
 6/(p+m),&0<p<6,\quad m\le6-p,\\
 1,&0<p<6,\quad m>6-p,
 \end{cases}\qquad
 \mu_a=\lambda-1.
\]
When $q>0$, $\mu_q=0$; when $q=0$, $\mu_q=p-6\ge0$. These formulas agree on the boundaries. In particular, a binding zero lower bound on a good need not have a strictly positive multiplier.
\paragraph{(d) Two incomes.} At $p=2$,
\[
 \begin{array}{c|ccc}
 m&a^*&q^*&\lambda\\\hline
 2&0&1&3/2\\
 8&4&2&1
 \end{array}
\]
With $m=2$, the no-borrowing/nonnegativity constraint on $a$ binds; the student would prefer to buy more tutoring at these prices but cannot finance it. For $m<4$, each additional unit of income buys $1/2$ more tutoring hour. With $m=8$, tutoring is at its unconstrained preferred level of 2 hours and additional income goes to $a$. Quasilinearity does not eliminate low-income corner effects.
\paragraph{(e) Indirect utility and Roy.} In the open region $0<p<6$, $m>6-p$, both goods are consumed and
\[
 v(p,m)=m-6+p+6\log(6/p),\qquad
 v_m=1,\qquad v_p=1-\frac6p.
\]
Hence $-v_p/v_m=6/p-1=q^*$. The derivative $v_m$ is the marginal value of income in the chosen utility units, equal to the budget multiplier. In this region one extra unit of income buys one more unit of $a$ and raises utility by 1. The multiplier is not invariant to rescaling the utility representation.

\clearpage\section*{A4. Short-answer key}
Award up to 5 points per answer. Respect the three-sentence limit: grade the argument, not literary style. Do not require students to use behavioural terminology not introduced in the course.
\subsection*{1. Paying to resist temptation}
\emph{Example answer (three sentences).} With fixed preferences and costless control over choices, removing an option for a fee cannot improve the best attainable outcome. Paying can nevertheless make sense if the app reduces attention costs or commits a future self whose choices differ from today's plans. The explanation must specify which feature is missing from the simple static model rather than treating every restriction as beneficial.

Marking: 2 points for the fixed-preference/costless-choice benchmark, 2 for a coherent attention or commitment explanation and 1 for explicitly identifying the changed assumption. A mere claim that ``everyone has different tastes'' is insufficient. If the explanation values the app's effect on concentration directly, it must make clear that the product then includes a different environment, not merely a smaller set of otherwise identical options.
\subsection*{2. Too cheap to be desirable}
\emph{Example answer (three sentences).} The shopper may value exclusivity or publicly recognised status, so a widely available bag need not be the same fully specified consumption alternative. Price may also provide information about quality. The statement alone does not establish incomplete or intransitive preferences; a meaningful comparison must hold relevant characteristics, information and social context fixed.

Marking: 2 points for a coherent characteristic or information explanation, 2 for stating what must be held fixed and 1 for avoiding an unsupported inference of irrationality. An unexplained label such as ``Veblen good'' is not a complete answer. Both price and availability change in the question, so do not infer a pure own-price effect from it.
\clearpage\section*{Part B · Training solutions}
\subsection*{B1. A preference checklist · 15 points}
For all three functions, representation gives completeness and transitivity. Each index is continuous, so weak comparisons survive limits.
\begin{center}\small\begin{tabular}{lccc}
Property & $u_1$ & $u_2$ & $u_3$\\\hline
Complete, transitive, continuous & Yes & Yes & Yes\\
Convex preferences & Yes & Yes & Yes\\
Strictly convex preferences & No & No & Yes\\
Strongly monotone & Yes & No & No\\
LNS everywhere on $X$ & Yes & Yes & No
\end{tabular}\end{center}
$u_1$ has half-space upper contours. Distinct indifferent bundles such as $(1,0)$ and $(0,2)$ remain indifferent when mixed, so strict convexity fails. Its positive coefficients imply strong monotonicity and LNS.

$u_2$ has rectangular upper contours. It is not strictly convex: $(2,1)$ and $(3,1)$, and their midpoint, all have utility 2. The same pair refutes strong monotonicity. Increasing both coordinates slightly strictly raises utility, establishing LNS at every bundle.

$u_3$ is strictly concave, hence represents strictly convex preferences: for distinct bundles the utility of an interior mixture is strictly above the corresponding utility average. Strong monotonicity fails since $(2,3)\gg(1,2)$ but has lower utility. LNS fails at the global bliss point $(1,2)$, so it is not a property holding everywhere (off that point, a small move toward it improves utility).

The requested plots through $(2,2)$ have levels 6, 2 and $-1$:
\begin{itemize}
\item $u_1=6$: segment $2x_1+x_2=6$ between $(0,6)$ and $(3,0)$; upper contour $2x_1+x_2\ge6$ in the orthant.
\item $u_2=2$: rays $x_1=2,x_2\ge1$ and $x_2=1,x_1\ge2$; upper contour $[2,\infty)\times[1,\infty)$.
\item $u_3=-1$: circle $(x_1-1)^2+(x_2-2)^2=1$, centred at $(1,2)$, tangent to $x_1=0$; upper contour is its closed disk.
\end{itemize}
Suggested subpart weights: (a) 3, (b) 8, (c) 4.
\subsection*{B2. Representations and group indices · 10 points}
(a) $3u-7$ is strictly increasing in $u$, and $u^2$ is strictly increasing on its attained range $[0,\infty)$, so both represent the same preferences. The truncated index does not: $(2,0)$ and $(3,0)$ have original utilities 2 and 3 but both truncate to 2.

(b) Level 2: line from $(0,1)$ to $(2,0)$. Level 4: line from $(0,2)$ to $(4,0)$. Under $3u-7$ they receive labels $-1$ and 5; under $u^2$ they receive labels 4 and 16.

(c) Use $v(t,c)=\sqrt{tc}$, $U(v,z)=(1+v)(1+z)$. Since $z\ge0$, $U$ is strictly increasing in $v$. At a common $z$, comparison reduces to comparing $v$, independently of that $z$. Separability does not specify a fixed group expenditure. Suggested weights: 4, 3, 3.
\clearpage
\subsection*{B3. Budget geometry · 10 points}
(a) $B=\{x\ge0:2x_1+3x_2\le18\}$; line $x_2=6-(2/3)x_1$, intercepts $(9,0),(0,6)$, slope $-2/3$.

(b) Costs are respectively 18, 21 and 15. The first and third are affordable; only the first exhausts the budget.

(c) With income 36 and the original prices, intercepts become $(18,0),(0,12)$ with unchanged slope. With $p_1=4$ at original income, intercepts are $(9/2,0),(0,6)$ and slope $-4/3$: an inward rotation around the vertical intercept.

(d) Scaling both prices and income by 2 leaves the original inequality unchanged after division by 2. If an optimum has slack $s=m-p\cdot x^*>0$, a ball of radius $s/(2\|p\|)$ around it stays affordable within the consumption domain. LNS supplies a better nearby bundle, contradicting optimality. This proof assumes an optimum exists. Suggested weights: 2, 2, 3, 3.
\clearpage
\subsection*{B4. Cobb--Douglas, one step at a time · 20 points}
(a) Original $u$ is continuous on the nonempty compact budget, so existence follows. Every axis bundle has utility zero; a strictly positive affordable bundle has positive utility. Thus every optimum has both quantities positive, and only then may the logarithm be used.

(b) $g=(1/3)\log x_1+(2/3)\log x_2$. For the interior optimum,
\[
 \mathcal L=g+\lambda_g(m-p_1x_1-p_2x_2),\qquad
 \frac1{3x_1}=\lambda_gp_1,\quad
 \frac2{3x_2}=\lambda_gp_2,\quad p\cdot x=m.
\]
(c) Summing expenditure equations gives $\lambda_g=1/m$ and
\[
 x_1^*=\frac{m}{3p_1},\qquad x_2^*=\frac{2m}{3p_2}.
\]
Shares are $1/3$ and $2/3$, and total expenditure is $m$. The log objective is strictly concave on the positive orthant, so the stationary solution is globally optimal among positive feasible bundles and unique; original boundary bundles are inferior by (a).

(d) Demand is $(2,8)$. At prices $(6,3)$ and income 36 it is still $(2,8)$.

(e) For the log index,
\[
 V_g=\tfrac13\log\frac{m}{3p_1}+\tfrac23\log\frac{2m}{3p_2},\qquad
 (V_g)_m=\frac1m,\quad (V_g)_{p_1}=-\frac1{3p_1},\quad (V_g)_{p_2}=-\frac2{3p_2}.
\]
Roy gives the demands above. (f) With $U^*=u(x^*)=e^{V_g}$, the original multiplier is $\lambda_u=U^*/m$, whereas $\lambda_g=1/m$. The log transformation divides marginal utilities and the multiplier by $U^*$; it preserves the ranking and chosen bundle. Suggested weights: 3, 4, 4, 2, 4, 3.
\subsection*{B5. Practising corners with Kuhn--Tucker · 20 points}
(a) Let $\mathcal L=u+\lambda(m-x_1-x_2)+\mu_1x_1+\mu_2x_2$. Stationarity is
\[
 \frac2{1+x_1}-\lambda+\mu_1=0,\qquad
 \frac1{1+x_2}-\lambda+\mu_2=0.
\]
Include $x_i\ge0$, $m-x_1-x_2\ge0$, $\lambda,\mu_i\ge0$ and $\lambda(m-x_1-x_2)=\mu_1x_1=\mu_2x_2=0$.

(b) With both goods positive, $\mu_1=\mu_2=0$, the budget binds, and
\[
 x_1=\frac{1+2m}{3},\qquad x_2=\frac{m-1}{3},\qquad\lambda=\frac3{m+2}.
\]
This is strictly interior iff $m>1$.

(c) With $x_2=0$ and budget exhausted, $x_1=m>0$, $\mu_1=0$, $\lambda=2/(1+m)$ and $\mu_2=\lambda-1=(1-m)/(1+m)$. Dual feasibility holds iff $m\le1$.

(d) The first regime $0<m\le1$ chooses $(m,0)$; for $m>1$, use the interior formula. Both agree at $m=1$. Concavity makes KKT sufficient, and strict concavity of the original utility in both quantities gives uniqueness.

(e) The table is
\[
 \begin{array}{c|ccccc}
 m&x_1&x_2&\lambda&\mu_1&\mu_2\\\hline
 1/2&1/2&0&4/3&0&1/3\\
 1&1&0&1&0&0\\
 4&3&1&1/2&0&0
 \end{array}
\]
At $m=1$ the zero-consumption constraint binds with zero multiplier. Suggested weights: 4, 4, 4, 4, 4.
\clearpage
\subsection*{B6. A kink without derivatives · 10 points}
(a) Attaining at least $k$ requires $x_1\ge k$ and $x_2\ge k/2$. Positive prices make the unique cheapest bundle $(k,k/2)$, with cost $(p_1+p_2/2)k$.

(b) The largest affordable target and demands are
\[
 k^*=\frac{2m}{2p_1+p_2},\quad
 x_1^*=\frac{2m}{2p_1+p_2},\quad x_2^*=\frac{m}{2p_1+p_2}.
\]
The cheapest implementation of the highest target costs exactly $m$. Any extra quantity would exceed the budget, so no other bundle attains that target affordably: uniqueness follows.

(c) At $(p_1,p_2,m)=(1,2,12)$, the bundle is $(6,3)$ and indirect utility is 6. (d) The index has a kink at $x_1=2x_2$, including the optimum. A differentiable MRS tangency calculation is inapplicable; the two target inequalities give the exact fixed proportions and cost. Suggested weights: 3, 3, 2, 2.
\subsection*{B7. Quasilinearity and the income boundary · 15 points}
(a) $a$ enters linearly with coefficient 1, so common feasible translations in $a$ preserve comparisons. The index is strictly increasing in $a$, implying exhaustion. Reduce to $\max_{0\le q\le m}\{m-q+5\log(1+q)\}$.

(b) The derivative is $-1+5/(1+q)$, with a zero at $q=4$; the second derivative is strictly negative. Thus $q^*=\min\{m,4\}$ and $a^*=\max\{m-4,0\}$.

(c) At $m=2$: $(a,q)=(0,2)$; at $m=6$: $(2,4)$. The graph of $q^*(m)$ follows the 45-degree line up to $m=4$ and is horizontal at 4 thereafter. The no-borrowing boundary on $a$ binds at low incomes.

(d) Indirect utility is
\[
 V(m)=\begin{cases}5\log(1+m),&0<m\le4,\\m-4+5\log5,&m>4.\end{cases}
\]
Both values equal $5\log5$ at 4. The left derivative is $5/(1+m)$, reaching 1 at 4; the right derivative is 1. Thus the value function is continuously differentiable across the regime change, although quantity demand has a kink.

(e) $\lambda=5/(1+m)$ for $m\le4$ and $\lambda=1$ for $m>4$. The nonnegativity multiplier is $\mu_a=\lambda-1$, while $\mu_q=0$ since $q^*>0$ at positive incomes. Below 4, extra income buys more tutoring/service quantity $q$; above 4 it buys only $a$. Suggested weights: 3, 3, 3, 3, 3.

\end{document}
